\documentclass[11pt,a4paper]{article}

\usepackage[margin=0.85in]{geometry}
\usepackage{parskip}
\setlength{\parskip}{0.45em}
\usepackage{hyperref}

\pagestyle{empty}

\begin{document}

\noindent
Alexei Vazquez\\
Nodes \& Links Ltd\\
Salisbury House, Station Road, Cambridge, CB1 2LA, UK\\
\texttt{alexei@nodeslinks.com}

\vspace{0.8em}
\noindent 6 August 2026

\vspace{1em}
\noindent Dear Dr Cerqueti,

\vspace{0.3em}

{\noindent\small
\textbf{Re: revised manuscript, submission ID dd2c9e03-9d4e-49f1-9513-ee3a977a7428}\\
\textbf{``Construction project dependency networks compress exponentially into
logical-depth levels whose mean distance to project completion is nearly
independent of project size''}\par}

\vspace{0.3em}

I am pleased to submit a revised version of the above manuscript for
consideration as an Article in \emph{Scientific Reports}. I am grateful to the
three reviewers for a careful and genuinely useful reading. Every point raised
has been addressed, and a detailed point-by-point response is enclosed
separately. The revision is substantial: the paper has been restructured, two of
its claims have been weakened, one has been settled by new analysis, and all
results are now reproducible from code and public data.

The statistical support the reviewers asked for is now in place: both size
relations are selected among four functional forms fitted under two error
structures on a common likelihood scale, with bootstrap confidence intervals,
residual diagnostics, and nested likelihood ratio tests of whether the three
companies share parameters. A descriptive statistics table makes clear that the
study rests on 1198 schedules rather than the handful one reviewer understandably
inferred. One relation explains only 27\% of the variance and does not
extrapolate to public benchmark data, so I no longer call it a scaling law. The
small-world claim, rightly called close to tautological as originally supported,
is now settled: I give a clustering coefficient appropriate to directed acyclic
graphs, build two degree-preserving null models, and show that the short
distances follow from the degree distributions of the coarse-grained network
rather than from any special arrangement of the long-range dependencies. The
results are replicated on the public PSPLIB benchmark; the engagement with the
design-structure-matrix, graph-drawing and network-symmetry literatures is
expanded; the explanatory schematic is now Figure 1; and the finer construction
is renamed from DFS to DPS to remove a clash with depth-first search. Following
your comment, the title has been replaced by a single 21-word sentence naming
the system studied and the two principal findings.

I would like to draw your attention to one matter directly. While rebuilding the
analysis for public release, I found a defect in the code used for the original
submission: a predecessor-pattern label was accumulated in floating point, which
silently merged distinct patterns in schedules deeper than 53 levels and
under-counted one reported quantity. I recomputed every number from the raw
schedules with exact arithmetic. All other results reproduced identically; the
affected exponent changes from 0.694 to 0.709 and one table column shifts
slightly, leaving the conclusions unchanged. The defect and its consequences are
described in full in the response, and the fix is public. I thought it important
to report this openly rather than substitute the corrected values silently.

The manuscript is original work, is not under consideration elsewhere, and I am
the sole author with no conflicts of interest to disclose. The implementation and
the scripts that generate every figure, table and statistical test are available
at \url{https://github.com/av2atgh/dus}. The construction schedules are
proprietary; the paper reports their summary statistics and shows the central
results are reproducible on public data.

Thank you for your time, and for the reviewers' effort, which has improved the
paper considerably.

\vspace{-0.4em}

\vspace{0.3em}
\noindent\begin{minipage}{\linewidth}
Sincerely,\\[1.4em]
Alexei Vazquez\\
Nodes \& Links Ltd
\end{minipage}

\end{document}
